\documentclass[11pt]{amsart}

\usepackage[margin=1in]{geometry}
\usepackage{amsmath,amssymb,mathtools}
\usepackage[T1]{fontenc}
\usepackage{lmodern}
\usepackage{microtype}
\usepackage[hidelinks]{hyperref}
\hypersetup{
  pdftitle={A 2-adic rigidity theorem for complementary divisors of Mersenne numbers},
  pdfauthor={Yuto Niwa},
  pdfsubject={An elementary rigidity theorem for complementary factors of prime-index Mersenne numbers},
  pdfkeywords={Mersenne numbers, 2-adic valuation, infinite descent, Lucas sequences}
}

\newtheorem{theorem}{Theorem}[section]
\newtheorem{lemma}[theorem]{Lemma}
\newtheorem{corollary}[theorem]{Corollary}
\theoremstyle{remark}
\newtheorem{remark}[theorem]{Remark}

\newcommand{\vTwo}{\nu_2}
\newcommand{\Mp}{M_p}

\title[A 2-adic rigidity theorem]{A 2-adic rigidity theorem for complementary divisors of Mersenne numbers}
\author{Yuto Niwa}
\address{Independent Researcher, Japan}
\email{yutoniwa0118@gmail.com}
\date{August 26, 2026}

\subjclass[2020]{Primary 11A51; Secondary 11B39, 11D61}
\keywords{Mersenne number, complementary divisors, 2-adic valuation, infinite descent, Lucas sequence}

\begin{document}

\begin{abstract}
Let $p$ be an odd prime and let $M_p=2^p-1$.  We prove that if
$M_p=DE$, where $D,E>1$, and
\[
 D\equiv 1\pmod{2^t},\qquad E\equiv -1\pmod{2^t}
\]
for an integer $t>p/4$, then necessarily
$(p,D,E)=(11,89,23)$.  Thus, apart from the factorization
$2^{11}-1=89\cdot23$, a complementary divisor pair of a prime-index
Mersenne number cannot lie simultaneously this deep in the residue
classes $1$ and $-1$ at the prime $2$.  The proof is elementary.  After
normalization, the factorization splits into two Diophantine branches.
One branch is impossible by infinite descent; the other descends to a
second-order Lucas-type recurrence, whose relevant power-of-two terms
are classified by congruences and a doubling identity.
\end{abstract}

\maketitle

\section{Introduction}

Write $M_n=2^n-1$ for the $n$th Mersenne number.  The prime divisors
and divisor statistics of Mersenne numbers have been investigated in,
among other works, \cite{Brillhart1964,Ehrman1967,Wagstaff1983,EKP1991,Cambraia2022}.
The present note concerns a local property of an arbitrary
\emph{complementary divisor pair}: a factorization $M_p=DE$ in which
one factor is close to $1$ and the other is close to $-1$ in the
$2$-adic metric.

Our main result is the following rigidity statement.

\begin{theorem}[2-adic rigidity]\label{thm:main}
Let $p$ be an odd prime, and suppose that
\[
 2^p-1=DE,\qquad D,E>1.
\]
If, for some integer $t>p/4$,
\[
 D\equiv 1\pmod{2^t},
 \qquad
 E\equiv -1\pmod{2^t},
\]
then
\[
 p=11,\qquad D=89,\qquad E=23.
\]
\end{theorem}

The condition is equivalently $t\ge \lfloor p/4\rfloor+1$.  A useful
valuation form follows immediately.

\begin{corollary}\label{cor:valuation}
Let $p>11$ be prime and let $2^p-1=DE$, where $D,E>1$,
$D\equiv1\pmod 8$, and $E\equiv-1\pmod 8$.  Then
\[
 \vTwo(D-1)=\vTwo(E+1)
 \le \left\lfloor\frac p4\right\rfloor.
\]
\end{corollary}

Primality of the exponent is essential.  Indeed, for every integer
$\tau\ge2$, putting $x=2^\tau$ gives
\[
 2^{3\tau}-1=(x^2+x+1)(x-1),
\]
and the two factors are congruent to $1$ and $-1$ modulo $2^\tau$,
respectively.  Here $\tau>(3\tau)/4$.  This family also explains why
the boundary case $p=3\tau$ must be treated explicitly below.

The proof has three stages.  First, the exact $2$-adic depths of the
two factors are shown to agree, and elementary size estimates force
that depth into a narrow interval.  Next, the normalized factorization
is reduced to two Diophantine systems, denoted N and P.  Finally, the
P-system is eliminated by infinite descent, while the N-system descends
to a Lucas-type recurrence.  Every step is included for completeness.

\section{Normalization}

We use $\vTwo(m)$ for the exponent of $2$ in a nonzero integer $m$.

\begin{lemma}\label{lem:equal-depth}
Under the hypotheses of Theorem~\ref{thm:main},
\[
 \vTwo(D-1)=\vTwo(E+1)=:\tau.
\]
Moreover,
\[
 3\tau<p<4\tau.
\]
\end{lemma}

\begin{proof}
Write
\[
 D=1+2^\alpha A_0,
 \qquad
 E=-1+2^\beta C_0,
\]
where $A_0$ and $C_0$ are positive and odd.  Expanding $DE+1=2^p$
gives
\begin{equation}\label{eq:initial-expansion}
 2^p=-2^\alpha A_0+2^\beta C_0
      +2^{\alpha+\beta}A_0C_0.
\end{equation}
If $\alpha<\beta$, the right-hand side has $2$-adic valuation exactly
$\alpha$, so $p=\alpha$.  But then $D>2^p$, whereas $D<DE=2^p-1$.
This is impossible.  The case $\beta<\alpha$ is symmetric.  Hence
$\alpha=\beta=\tau$.

Set $x=2^\tau$ and write
\[
 D=xA+1,
 \qquad
 E=xC-1,
\]
with $A,C$ positive and odd.  Since
\[
 2^p-1=DE\ge(x+1)(x-1)=x^2-1,
\]
we have $p\ge2\tau$.  As $p$ is odd, in fact $p>2\tau$.
Expanding the factorization now yields
\[
 x^2AC+x(C-A)=2^p.
\]
Consequently $x\mid C-A$; write
\[
 C-A=jx.
\]
After division by $x^2$ we obtain
\begin{equation}\label{eq:ACj}
 AC+j=2^{p-2\tau}.
\end{equation}

Suppose first that $p<3\tau$.  Then the right-hand side of
\eqref{eq:ACj} is smaller than $x$.  If $j>0$, then $C=A+jx$ and
$AC+j>x$.  If $j=0$, then $A=C$ and an odd square would equal a
power of $2$ greater than $1$.  If $j=-k<0$, then $A=C+kx$ and
\[
 AC+j=C^2+k(xC-1)\ge x.
\]
Each case is a contradiction.

It remains to exclude $p=3\tau$.  The same three cases show that
equality in \eqref{eq:ACj} can occur only when $j=-1$ and $C=1$.
Then $A=x+1$ and
\[
 DE=(x^2+x+1)(x-1)=x^3-1.
\]
Since $p=3\tau$ is prime, necessarily $(p,\tau)=(3,1)$; but then
$E=x-1=1$, contrary to the hypotheses.  Thus $p>3\tau$.  Finally,
$\tau\ge t>p/4$ gives $p<4\tau$.
\end{proof}

Define
\begin{equation}\label{eq:Hy-def}
 H=2^{p-3\tau},
 \qquad
 y=2^{4\tau-p}.
\end{equation}
Lemma~\ref{lem:equal-depth} implies
\begin{equation}\label{eq:Hy-properties}
 x=Hy,
 \qquad H\ge2,
 \qquad y\ge2,
 \qquad H<x,
\end{equation}
and \eqref{eq:ACj} becomes
\begin{equation}\label{eq:xH}
 AC+j=xH.
\end{equation}
Both $H$ and $y$ are powers of $2$.

\section{Reduction to two Diophantine branches}

\begin{lemma}\label{lem:two-branches}
There exist positive integers $H,a,q,r$ such that $q$ and $r$ are
odd and $\gcd(a,q)=1$, and exactly one of the following systems holds:
\begin{align}
 H&=qr+a, & r^2&=aHy+q; \tag{N}\label{eq:N}\\[2mm]
 H&=qr+a, & r^2&=aHy-q. \tag{P}\label{eq:P}
\end{align}
\end{lemma}

\begin{proof}
First we show that $A,C<x^2$.  Since $H$ is a proper power-of-two
divisor of $x$, we have $H\le x/2$.  If $A\ge x^2$, then
\[
 DE\ge(x^3+1)(x-1)
    >\frac{x^4}{2}-1
    \ge x^3H-1
    =2^p-1,
\]
a contradiction.  If $C\ge x^2$, the analogous estimate
\[
 DE\ge(x+1)(x^3-1)>\frac{x^4}{2}-1
\]
gives the same contradiction.

Because $A\equiv C\pmod x$, there are integers $r,s,u$ such that
\[
 A=r+sx,
 \qquad
 C=r+ux,
\]
where $1\le r<x$, $r$ is odd, and $0\le s,u<x$.  Since
$j=(C-A)/x=u-s$, substitution in \eqref{eq:xH} gives
\[
 r^2+r(s+u)x+sux^2+u-s=xH.
\]
It follows that
\[
 a:=\frac{r^2+u-s}{x}
\]
is an integer.  The numerator is a multiple of $x$ strictly between
$-x$ and $x^2$, so $0\le a<x$.  Dividing the preceding identity by
$x$ yields
\begin{equation}\label{eq:Hrsu}
 H=a+r(s+u)+sux.
\end{equation}
Because $H<x$ and all terms on the right are nonnegative, $su=0$.
The case $s=u=0$ would give $x\mid r^2$, impossible because $x$ is
even and $r$ is odd.

If $u=0$ and $s=q>0$, then \eqref{eq:Hrsu} and the definition of
$a$ give the N-system.  If $s=0$ and $u=q>0$, they give the
P-system.  In either case $q$ is odd, because $x$ divides $r^2-q$
or $r^2+q$.  Also $a>0$: if $a=0$, then $H=qr$ would be odd,
contrary to $H\ge2$ being a power of $2$.  Finally, every common
divisor of $a$ and $q$ divides $H=a+qr$.  Such a divisor is odd,
while $H$ is a power of $2$; hence $\gcd(a,q)=1$.
\end{proof}

\section{Infinite descent}

We first dispose of the P-branch.  The following lemma is deliberately
stated without a coprimality assumption; this makes the descent closed
under iteration.

\begin{lemma}\label{lem:P}
Let $y\ge2$ be even and let $q,r$ be positive odd integers.  There are
no positive integers $H,a$ satisfying
\[
 H=qr+a,
 \qquad
 r^2=aHy-q.
\]
\end{lemma}

\begin{proof}
Assume that a solution exists.  Since $H=qr+a>r$, eliminating $q$
gives
\begin{equation}\label{eq:P-eliminate}
 H(ayr-1)=r^3-a.
\end{equation}
Comparison with $H>r$ gives
\[
 r^2(r-ay)+(r-a)>0.
\]
Because $r$ is odd and $ay$ is even, equality $r=ay$ is impossible.
If $r<ay$, the left-hand side is at most $-r^2+r-a<0$.
Therefore
\begin{equation}\label{eq:rgtay}
 r>ay.
\end{equation}
In particular $r\ge3$.  From \eqref{eq:P-eliminate} we then obtain
\[
 H=\frac{r^3-a}{ayr-1}
   <\frac{r^3}{2r-1}
   <r^2,
\]
and hence
\begin{equation}\label{eq:qltr}
 q=\frac{H-a}{r}<r.
\end{equation}

Set
\[
 c=r-aqy.
\]
A direct calculation using the P-system gives
\begin{equation}\label{eq:Pcr}
 cr=a^2y-q.
\end{equation}
The integer $c$ is odd.  If $c\le0$, then $-c\ge1$, and
\eqref{eq:Pcr} together with \eqref{eq:qltr} gives
\[
 r\le(-c)r=q-a^2y<q<r,
\]
a contradiction.  Thus $c>0$.  If $c\ge a$, then
\eqref{eq:Pcr} gives $r<ay$, contradicting \eqref{eq:rgtay}.
Consequently
\begin{equation}\label{eq:P-c-range}
 0<c<a.
\end{equation}

Put $b=a-qc$.  Substituting $r=aqy+c$ in \eqref{eq:Pcr} gives
\begin{equation}\label{eq:Pdesc}
 c^2=aby-q,
 \qquad
 a=qc+b.
\end{equation}
The first identity shows that $b>0$.  If $b\ge c$, then
\[
 c^2+q=aby\ge2(q+1)c^2>c^2+q,
\]
which is impossible.  Hence $0<b<c$.

Equations \eqref{eq:Pdesc} say that $(a,c,b)$ is another P-solution:
its variables $(H,r,a)$ are replaced by $(a,c,b)$, while $q$ and $y$
are unchanged.  All required positivity and parity conditions are
preserved, and the first component decreases because $a<H=qr+a$.
Iteration would therefore produce an infinite strictly decreasing
sequence of positive integers, a contradiction.
\end{proof}

We turn to the N-branch.

\begin{lemma}\label{lem:N}
Let $y\ge2$ be even and let $q,r$ be positive odd integers.  Suppose
that positive integers $H,a$ satisfy
\[
 H=qr+a,
 \qquad
 r^2=aHy+q,
 \qquad
 \gcd(a,q)=1.
\]
Then $q=1$.  Repeated descent terminates at
\[
 H=y+2,
 \qquad
 r=y+1,
 \qquad
 a=1.
\]
\end{lemma}

\begin{proof}
Set
\[
 c=r-aqy.
\]
The N-system gives
\begin{equation}\label{eq:Ncr}
 cr=a^2y+q,
\end{equation}
so $c>0$.  We claim that $c\le a$.  If $c>a$, then
\[
 r=aqy+c>aqy+a,
\]
whereas \eqref{eq:Ncr} gives
\[
 r=\frac{a^2y+q}{c}<ay+\frac qa.
\]
Combining the two inequalities yields
\begin{equation}\label{eq:N-impossible-c-large}
 a^2y(q-1)+a^2<q.
\end{equation}
For $q=1$ this says $a^2<1$.  For the remaining odd values $q\ge3$,
the left-hand side is at least $y(q-1)+1\ge2q-1>q$.  Thus
\eqref{eq:N-impossible-c-large} is always impossible.

If $c=a$, then \eqref{eq:Ncr} gives $q=a(r-ay)$.  Coprimality forces
$a=1$.  Now $c=r-qy=1$ and $q=r-y$, so
\[
 (q-1)(y-1)=0.
\]
Since $y\ge2$, we obtain $q=1$, $r=y+1$, and $H=y+2$.

It remains to consider $0<c<a$.  Put $b=a-qc$.  Substitution in
\eqref{eq:Ncr} gives
\begin{equation}\label{eq:Ndesc}
 c^2=aby+q,
 \qquad
 a=qc+b.
\end{equation}
We show that $b>0$.  If $b<0$, set $e=-b=qc-a>0$.  Then
\[
 q-c^2=aye>0.
\]
As $1\le e\le qc-1$, the elementary bound
$e(qc-e)\ge qc-1$ gives
\[
 q-c^2=aye\ge2(qc-1).
\]
But
\[
 2(qc-1)-(q-c^2)=q(2c-1)+c^2-2>0:
\]
for $c=1$ this is $q-1>0$ because $q>c^2$, and for $c\ge3$ it is
immediate.  This is a contradiction.  If $b=0$, then $q=c^2$ and
$a=qc$.  Coprimality would force $q=c=a=1$, contrary to $c<a$.
Hence $b>0$.

Since $c<a$ and $y\ge2$, we have $ay>c$.  The first identity in
\eqref{eq:Ndesc} therefore implies $b<c$.  Moreover,
\[
 \gcd(b,q)=\gcd(a-qc,q)=1.
\]
Thus $(a,c,b)$ is another N-solution with the same $q$ and $y$, and
its first component $a$ is smaller than $H$.  Repeated descent must
eventually reach the case $c=a$, which, as shown above, forces $q=1$
and gives the stated terminal solution.
\end{proof}

\section{The terminal recurrence}

It remains to determine which N-solutions with $q=1$ can have $H$ a
power of $2$.

\begin{lemma}\label{lem:recurrence}
Fix an even integer $y\ge2$.  Define
\[
 \binom{H_0}{r_0}=\binom11,
 \qquad
 \binom{H_{m+1}}{r_{m+1}}
 =
 \begin{pmatrix}y+1&1\\ y&1\end{pmatrix}
 \binom{H_m}{r_m}.
\]
Every positive N-solution with $q=1$ occurs uniquely among the pairs
$(H_m,r_m)$ with $m\ge1$.  In particular,
\begin{equation}\label{eq:H-recurrence}
 H_0=1,
 \qquad H_1=y+2,
 \qquad H_{m+2}=(y+2)H_{m+1}-H_m.
\end{equation}
If $y$ and some $H_m>1$ are both powers of $2$, then $y=2$ and
$m=1$.
\end{lemma}

\begin{proof}
For $q=1$, the inverse of the descent in Lemma~\ref{lem:N} sends a
smaller pair $(H',r')$ to
\[
 r=yH'+r',
 \qquad
 H=(y+1)H'+r'.
\]
The terminal solution $(y+2,y+1)$ is the image of $(1,1)$, so repeated
inversion gives the displayed matrix recurrence.  Conversely, the
relation
\[
 r^2-yH(H-r)=1
\]
is preserved by this transformation; hence every pair with $m\ge1$
is an N-solution with $a=H_m-r_m>0$.  Lemma~\ref{lem:N} proves that
every positive N-solution descends uniquely to the terminal one, so the
list is exhaustive and has no repetitions.

The matrix has trace $y+2$ and determinant $1$, which gives the scalar
recurrence \eqref{eq:H-recurrence}.  Modulo $2$ it gives
$H_{m+2}\equiv H_m$, so precisely the odd-indexed $H_m$ are even.
Thus, if $H_m>1$ is a power of $2$, then $m$ is odd.  Modulo $y+2$,
the same recurrence gives $H_{m+2}\equiv-H_m$.  Since
$H_1=y+2$, every odd-indexed term is divisible by $y+2$.  Therefore
$y+2$ is itself a power of $2$.  As $y$ is a positive power of $2$,
this forces
\begin{equation}\label{eq:y-two}
 y=2.
\end{equation}

For $y=2$, let
\[
 U_0=0,
 \quad U_1=1,
 \quad U_{k+2}=4U_{k+1}-U_k.
\]
Then $H_m=U_{m+1}$.  Define the companion sequence
\[
 V_0=2,
 \quad V_1=4,
 \quad V_{k+2}=4V_{k+1}-V_k.
\]
If $\lambda$ and $\lambda^{-1}$ are the roots of
$X^2-4X+1$, then
\[
 U_k=\frac{\lambda^k-\lambda^{-k}}{\lambda-\lambda^{-1}},
 \qquad
 V_k=\lambda^k+\lambda^{-k}.
\]
Consequently
\begin{equation}\label{eq:doubling}
 U_{2k}=U_kV_k.
\end{equation}

An induction in the recurrence modulo $8$ gives
\[
 V_k\equiv
 \begin{cases}
 2 & (k\equiv0\pmod4),\\
 4 & (k\equiv1\pmod4),\\
 6 & (k\equiv2\pmod4),\\
 4 & (k\equiv3\pmod4).
 \end{cases}
 \pmod8.
\]
Also $V_2=14$, and the positive sequence is strictly increasing from
that point onward.  Hence every $V_k$ with $k\ge2$ has an odd divisor
greater than $1$.  By \eqref{eq:doubling}, $U_{2k}$ cannot be a power
of $2$ for $k\ge2$.

Finally, $m$ is odd, so write $m+1=2k$.  The preceding paragraph
forces $k=1$, and therefore $m=1$.
\end{proof}

\section{Proofs of the main results}

\begin{proof}[Proof of Theorem~\ref{thm:main}]
Lemmas~\ref{lem:equal-depth} and \ref{lem:two-branches} reduce the
factorization to the N-system or the P-system.  Lemma~\ref{lem:P}
excludes the P-system.  In the N-system, Lemma~\ref{lem:N} gives
$q=1$, while Lemma~\ref{lem:recurrence}, applied to the powers of two
$y$ and $H$, gives
\[
 y=2,
 \qquad H=4.
\]
Since $x=Hy$, it follows that $x=8$ and $\tau=3$.  The definition
$H=2^{p-3\tau}$ then gives $p=11$.

The relevant N-solution is the terminal one:
\[
 r=y+1=3,
 \qquad a=q=1.
\]
In the notation of Lemma~\ref{lem:two-branches}, the N-branch arose
from $u=0$ and $s=q$, so
\[
 A=r+qx=11,
 \qquad C=r=3.
\]
Therefore
\[
 D=xA+1=89,
 \qquad E=xC-1=23.
\]
Conversely, $2^{11}-1=89\cdot23$, with
$89\equiv1\pmod8$ and $23\equiv-1\pmod8$, so the exceptional case
does satisfy the hypotheses.
\end{proof}

\begin{proof}[Proof of Corollary~\ref{cor:valuation}]
The first part of Lemma~\ref{lem:equal-depth} gives equality of the two
valuations.  If their common value were greater than $p/4$, then
Theorem~\ref{thm:main} would force $p=11$.  Since the common value is
an integer, the stated bound follows.
\end{proof}

\section{Concluding remarks}

Theorem~\ref{thm:main} applies to arbitrary complementary factors;
neither factor is assumed prime, and no hypothesis is made on the
number of prime factors of $M_p$.  The exceptional factorization
$2^{11}-1=23\cdot89$ has a structural explanation in the proof: it
corresponds to the first nontrivial term $H_1=4$ of the terminal
Lucas-type recurrence, whereas every later relevant term has an odd
divisor.

The theorem does not imply that there are infinitely many Mersenne
primes.  It only rules out one particularly deep form of simultaneous
$2$-adic cancellation in a composite prime-index Mersenne number.

\section*{Formal verification}
The main theorem has been formalized in Lean 4.25.2 with mathlib
v4.25.2. The ancillary file \texttt{MersenneRigidity.lean} contains the
proof of \texttt{MersenneRigidity.mersenne\_rigidity}.
The Lean proof was generated using ChatGPT.

\section*{Declarations}

\noindent\textbf{Competing interests.}
The author reports no competing interests.

\smallskip
\noindent\textbf{Generative AI use.}
ChatGPT generated nearly all of the manuscript text, including the written
proof. It was also used for proof exploration, bounded counterexample
searches, and literature searches. The main theorem was separately checked
by the Lean kernel. The human author is responsible for the manuscript,
including the statement of the theorem, the exposition, and the references.

\begin{thebibliography}{99}

\bibitem{Brillhart1964}
J.~Brillhart,
On the factors of certain Mersenne numbers. II,
\emph{Math. Comp.} \textbf{18} (1964), 87--92.
\href{https://doi.org/10.1090/S0025-5718-1964-0159776-X}{doi:10.1090/S0025-5718-1964-0159776-X}

\bibitem{Cambraia2022}
A.~Cambraia Jr., M.~P. Knapp, A.~Lemos, B.~K. Moriya, and
P.~H.~A. Rodrigues,
On prime factors of Mersenne numbers,
\emph{Palestine J. Math.} \textbf{11} (2022), no.~2, 449--456.
\href{https://pjm.ppu.edu/paper/1088-prime-factors-mersenne-numbers}{journal page}

\bibitem{Ehrman1967}
J.~R. Ehrman,
The number of prime divisors of certain Mersenne numbers,
\emph{Math. Comp.} \textbf{21} (1967), 700--704.
\href{https://doi.org/10.1090/S0025-5718-1967-0223320-1}{doi:10.1090/S0025-5718-1967-0223320-1}

\bibitem{EKP1991}
P.~Erd{\H{o}}s, P.~Kiss, and C.~Pomerance,
On prime divisors of Mersenne numbers,
\emph{Acta Arith.} \textbf{57} (1991), 267--281.
\href{https://doi.org/10.4064/aa-57-3-267-281}{doi:10.4064/aa-57-3-267-281}

\bibitem{Wagstaff1983}
S.~S. Wagstaff, Jr.,
Divisors of Mersenne numbers,
\emph{Math. Comp.} \textbf{40} (1983), 385--397.
\href{https://doi.org/10.1090/S0025-5718-1983-0679454-X}{doi:10.1090/S0025-5718-1983-0679454-X}

\end{thebibliography}

\end{document}
